% Part: first-order-logic
% Chapter: introduction
% Section: first-order-logic

\documentclass[../../../include/open-logic-section]{subfiles}

\begin{document}

\olfileid{fol}{int}{fol}

\olsection{Logika Sepisanan}

Sampeyan mbokmenawa wis wanuh karo logika sepisanan saka pambuka
pisanan marang logika formal.\footnote{Sejatine, kita kurang luwih
nganggep yen sampeyan pancen wis wanuh!{} Yen durung, sampeyan bisa
nyemak maneh buku ajar kang luwih dhasar, kayata \emph{forall x}
\citep{Magnus2021}.} Sampeyan bisa uga mangerteni logika iki kanthi
jeneng ``logika kuantifikasional'' utawa ``logika predikat.'' Kaping
pisan, logika sepisanan iku basa formal. Tegese, logika iki nduweni
tetembungan tartamtu, lan ekspresine awujud rerangken simbol saka
tetembungan kasebut. Nanging, ora saben rerangken simbol diparengake.
Ana maneka jinis ekspresi kang diparengake: term, !!{formula}s, lan
!!{sentence}s. Kawigaten kita utamane tumuju marang !!{sentence}s ing
logika sepisanan: ukara-ukara kasebut menehi padanan formal kanggo
ukara basa Inggris, lan kita bisa ngajokake pitakon-pitakon kang
lumrahe narik kawigatene ahli logika. Tuladhane:
\begin{itemize}
    \item Apa $!B$ ndherek saka~$!A$ kanthi logis?
    \item Apa $!A$ bener kanthi logis, salah kanthi logis, utawa
kontingen?
    \item Apa $!A$ lan $!B$ ekuivalen?
\end{itemize}

Pitakon-pitakon iki utamane babagan ``makna'' !!{sentence}s ing logika
sepisanan. Tuladhane, ahli filsafat bakal nganalisis pitakon apa $!B$
ndherek kanthi logis saka~$!A$ kanthi takon: apa ana kasus nalika $!A$
bener nanging~$!B$ salah ($!B$ ora ndherek saka~$!A$), utawa apa saben
kasus kang ndadekake $!A$ bener uga ndadekake~$!B$ bener ($!B$ ndherek
saka~$!A$)? Nanging, kita durung diwenehi katrangan apa kang diarani
``kasus''---iku tugase \emph{semantik}. Semantik logika sepisanan
menehi model matematis kang trep kanggo gagasan intuisi ahli filsafat
ngenani ``kasus,'' lan uga---iki wigati---ngenani apa tegese
!!a{sentence}~$!A$ \emph{bener ing} sawijining kasus. Model matematis
kang trep iki, kang bakal kita wangun, diarani !!a{structure}. Relasi
kang njlentrehake ``bener ing'' kanthi trep diarani relasi
\emph{nyukupi}. Mula, kang bakal kita temtokake yaiku ``$!A$
dicukupi ing~$\Struct{M}$'' (ing simbol: $\Sat{M}{!A}$) kanggo
!!{sentence}s~$!A$ lan !!{structure}s~$\Struct{M}$. Sawise iki rampung,
kita uga bisa menehi definisi kang trep kanggo istilah-istilah semantis
liyane, kayata ``ndherek saka'' utawa ``bener kanthi logis.''
Definisi-definisi kasebut bakal ndadekake kita bisa netepake, kanthi
presisi matematis maneh, apa, upamane,
$\lforall[x][(!A(x) \lif !B(x))],
\lexists[x][!A(x)] \Entails \lexists[x][!B(x)]$. Cathetan koreksi
sumber OLPL-056: sumber Inggris ora nutup argumen kuantor universal
sawise B(x), lan kurung kasebut kliru dumunung sawise kesimpulan,
mula cakupan formula premis rusak; kurung panutup dibalekake adhedhasar
rong pangulangan kang bener ing ngisor iki.
Jawabane mesthi ``ya.'' Yen sampeyan wis sinau nyimbolake ukara basa
Inggris ing logika sepisanan, sampeyan bakal ngenali iki minangka,
upamane, simbolisasi saka pratelan ``Kabeh semut iku serangga, ana
semut, mula ana serangga.'' Iki cetha minangka argumen kang valid,
mula model matematis kita kanggo ``ndherek saka'' ing basa formal kudu
menehi jawaban kang padha.

Topik liyane kang mbokmenawa isih sampeyan eling saka pambuka pisanan
marang logika formal yaiku anane \emph{!!{derivation}s}. Yen sampeyan
wis tau melu kuliah logika formal dhasar, pamulang mesthine wis menehi
latihan nggoleki !!{derivation}s kasebut, bisa uga malah !!a{derivation}
kang nuduhake yen akibat ing ndhuwur lumaku. Ana maneka cara kanggo
menehi !!{derivation}s: sampeyan bisa uga wis nggunakake cara kang
diarani ``deduksi natural'' utawa ``wit kabeneran,'' nanging isih ana
akeh cara liyane. Tujuane sistem !!{derivation} yaiku nyedhiyakake
piranti kanggo mangsuli pitakon-pitakone ahli logika ing ndhuwur:
upamane, !!{derivation} deduksi natural kang nduweni
$\lforall[x][(!A(x) \lif !B(x))]$ lan $\lexists[x][!A(x)]$ minangka
premis sarta $\lexists[x][!B(x)]$ minangka kesimpulan (larik pungkasan)
\emph{mesthekake} yen $\lexists[x][!B(x)]$ ndherek kanthi logis saka
$\lforall[x][(!A(x) \lif !B(x))]$ lan $\lexists[x][!A(x)]$. Cathetan
koreksi sumber OLPL-057: formula premis pisanan ing ukara Inggris iki
kelangan kurung panutup lan kurung kasebut kliru dipindhah menyang
formula kesimpulan; rong formula iki dibalekake menyang cakupane dhewe.

Nanging, apa sebabe? Yen mung dideleng saklumrahe, sistem
!!{derivation} kaya ora ana gandhengane karo semantik: menehi
!!{derivation} formal mung mbutuhake panyusunan simbol kanthi cara-cara
tartamtu kang manut aturan; sistem kasebut ora nyebut ``kasus'' utawa
``bener ing'' babar pisan. Gandhengane sistem !!{derivation} karo
semantik kudu ditetepake liwat panliten metalogis. Kang dibutuhake yaiku
bukti matematis, upamane, yen !!{derivation} formal kanggo
$\lexists[x][!B(x)]$ saka premis $\lforall[x][(!A(x) \lif !B(x))]$ lan
$\lexists[x][!A(x)]$ bisa digawe yen, lan mung yen,
$\lforall[x][(!A(x) \lif !B(x))]$ lan $\lexists[x][!A(x)]$ bebarengan
nduweni akibat~$\lexists[x][!B(x)]$. Cathetan koreksi sumber OLPL-058:
pratelan Inggris pungkasan ngulang pola kurung kang padha, yaiku kurung
panutup universal ilang lan dumunung keluwih sawise kesimpulan; saben
kurung dibalekake menyang formula kang trep. Nanging, sadurunge iki
bisa ditindakake, akeh pakaryan tliti kudu dirampungake supaya definisi
sintaksis lan semantik bener.

\end{document}
